%% ---------------------------------------------------------------------------
%% intro.tex
%%
%% Introduction
%%
%% $Id: intro.tex 1477 2010-07-28 21:34:43Z palvarado $
%% ---------------------------------------------------------------------------

\chapter{Introducción}
\label{chp:intro}

Agregar texto introductorio al capítulo

\section{Antecedentes del proyecto}

Agregar texto introductorio a la sección

\subsection{Descripción de la organización}

RECOMENDACIÓN: En la descripción de la organización no olvidar las referencias.

\subsection{Descripción del área de conocimiento del proyecto}

\textbf{Desarrollo de algoritmos en sistemas empotrados}

\subsection{Trabajos similares encontrados}

RECOMENDACIONES:

\begin{itemize}
    \item Es posible que no exista un trabajo idéntico al que está realizando por lo que puede buscar variantes. La idea es mostrar aproximaciones para un problema similar. Debe indicar detalles de implementación para encontrar alguna similitud con su trabajo.
    \item Algunos proyectos pueden presentar confidencialidad dentro de la empresa, por lo que pueden citar que se desarrolló un sistema con X y Y características (sin detallar nombres). Adicionalmente, puede aplicar el punto anterior.

\end{itemize}


\section{Planteamiento del problema}

Agregar texto introductorio a la sección

\subsection{Contexto del problema}

\subsection{Justificación del problema}

\subsection{Enunciado del problema}

no colocarlo en forma de solución. Lo que va a desarrollar en su proyecto no es el problema en sí del proyecto

\section{Objetivos del proyecto}

Agregar texto introductorio a la sección

\subsection{Objetivo general}

El objetivo general va a establecer cómo se va a resolver el problema que estableció anteriormente.

\subsection{Objetivos específicos}

numerados

Para los objetivos recuerde que debe contestar las preguntas qué, cómo y para.

Los objetivos no son actividades. Deben colocarse entre 3 y 4 específicos.

Los objetivos deben tener un verbo que termine en -ar, -er, -ir. Basarse en la taxonomía de Bloom.

Asegurarse de que cada objetivo específico cumpla con las siguientes características:
i.	Tener un entregable(s) definido.
ii.	Establecer un procedimiento para desarrollarlo.
iii.	Se pueda evidenciar el cumplimiento (resultados).


Para los entregables, asegurarse que se pueda evidenciar que se lograron.


\section{Alcances, entregables y limitaciones del proyecto}


%%% Local Variables: 
%%% mode: latex
%%% TeX-master: "main"
%%% End: 
